\section{Methodology and Implementation}
\label{sec:methodology}

This section describes the five datasets used, the data pipeline, the three
recurrent architectures, and the training, model-selection and evaluation procedures.

The implementation separates data preparation, recurrent models, training and
evaluation into shared Python modules. NumPy and pandas provide the data
pipeline, and PyTorch implements the networks and automatic differentiation. All
three architectures use the same eligible examples, scaling rules, loss function
and evaluation procedure. Their feedback connections and validation-selected
hidden sizes differ.

\subsection{Dataset Overview}
The following five datasets are used throughout the report. These short names
identify the forecasting tasks; Section~\ref{sec:datasets} provides their sources,
plots and more detailed discussion of their temporal properties.


\noindent\textbf{Air passengers:} Monthly international airline passenger totals,
  measured in thousands. The task is to predict the next month's total from
  previous totals in a series with growth and annual seasonality.

\noindent\textbf{Minimum temperatures:} Daily minimum temperatures in Melbourne,
  measured in degrees Celsius. The task is to predict the next day's minimum
  from past temperatures, which reflect annual seasonality and daily variation.

\noindent\textbf{Smart home:} Environmental sensor measurements from a monitored house,
  recorded at 15-minute intervals. Indoor and outdoor sensor histories are used
  to predict dining-room and bedroom temperatures at the next interval.

\noindent\textbf{Sunspots:} Monthly mean total sunspot numbers, used as an index of solar
  activity. The task is to predict the next month's value from previous values
  in a series with approximately 11-year cycles.

\noindent\textbf{AR(2):} A synthetic second-order autoregressive series in which each
  value depends on its two predecessors and added random noise. It provides a
  controlled stationary forecasting task with known generating dynamics.

\subsection{Forecasting Task and Data Preparation}
Let $\mathbf{x}_t\in\mathbb{R}^{d}$ contain the input features available at
forecast origin $t$, and let $\mathbf{y}_{t+q}\in\mathbb{R}^{m}$ contain the
targets $q$ steps ahead. An example consists of
\begin{equation}
 \mathbf{X}_t=(\mathbf{x}_{t-L+1},\ldots,\mathbf{x}_t),
 \qquad \text{target}=\mathbf{y}_{t+q},
\end{equation}
where $L$ is the lookback. A minibatch is a group of $B$ training examples
processed together to compute one weight update. Its input array has shape $B\times L\times d$, and its
predictions have shape $B\times m$. The pipeline supports configurable inputs,
targets and horizons, but every reported experiment uses $q=1$. Passenger,
minimum-temperature, sunspot and AR(2) experiments use the target's own history
($d=m=1$).
Smart-home experiments use dining-room and bedroom temperature,
outdoor temperature and outdoor humidity ($d=4$), predicting both indoor
temperatures jointly ($m=2$).

Dataset-specific scripts prepare regular time indices before window construction.
Passenger values require no imputation. The minimum-temperature cleaner removes
question-mark prefixes from three numeric entries, restores two missing dates
and forward-fills them. Smart-home recordings are processed separately: two
missing 15-minute records are inserted and forward-filled, but a long gap
between recordings is preserved. Observation flags are assigned before filling so
that inserted values can serve as input history without becoming ground-truth
targets. The sunspot cleaner preserves the supplied target series, including a
historically interpolated February 1824 value, while retaining quality metadata
separately from model inputs. AR(2) samples are generated directly with no
missing values; the generator settings are given in Section~\ref{sec:datasets}.

The shared loader validates timestamps and selected numeric features and assigns
a segment identifier to each uninterrupted recording. Any example whose
input-to-target interval crosses a segment boundary is rejected. Examples with no
observed target are also rejected. Validation and test windows may use preceding training
observations as context, since those observations are available at the forecast
origin.

Each input and target feature is standardized separately:
\begin{equation}
 z_{t,j}=\frac{x_{t,j}-\mu_j}{s_j}.
\end{equation}
The mean $\mu_j$ and population standard deviation $s_j$ are fitted using
observed rows in the current training interval only, before overlapping windows
are constructed. Thus repeated appearances of a row in different windows do not
change its scaling weight. Constant features use $s_j=1$. The same fitted
transformation is applied to the corresponding validation or test interval. No
differencing, logarithmic transformation or seasonal adjustment is used, and
predictions are inverse-transformed before reporting errors in original units.

\subsection{Recurrent Network Implementations}
Each network has one hidden layer of $H$ tanh units and a linear output layer.
Within an input window, let $k$ index the sequence position,
$\mathbf{h}_k\in\mathbb{R}^{H}$ the hidden activation, and
$\mathbf{o}_k\in\mathbb{R}^{m}$ the output. The Elman network uses hidden-state
feedback:
\begin{equation}
 \mathbf{h}_k=\tanh(W_x\mathbf{x}_k+W_h\mathbf{h}_{k-1}+\mathbf{b}_h).
\end{equation}
The Jordan network uses output feedback:
\begin{equation}
 \mathbf{h}_k=\tanh(W_x\mathbf{x}_k+W_o\mathbf{o}_{k-1}+\mathbf{b}_h).
\end{equation}
The multi-recurrent network combines both paths:
\begin{equation}
 \mathbf{h}_k=\tanh(W_x\mathbf{x}_k+W_h\mathbf{h}_{k-1}
                         +W_o\mathbf{o}_{k-1}+\mathbf{b}_h).
\end{equation}
All three compute
\begin{equation}
 \mathbf{o}_k=W_y\mathbf{h}_k+\mathbf{b}_y,
 \qquad \hat{\mathbf{y}}_{t+1}=\mathbf{o}_L.
\end{equation}
Here $W_x$ is $H\times d$, $W_h$ is $H\times H$, $W_o$ is $H\times m$, and
$W_y$ is $m\times H$. Each implementation has one hidden bias and one output
bias; feedback projections have no additional biases. The output layer is linear
because the task is regression on standardized real-valued targets.

Hidden and output contexts are initialized to zero independently for each window.
Context stores a one-step copy; no separate leaky context self-connection is
used. Jordan and multi-recurrent feedback always uses the network's own
intermediate output, not a supplied ground-truth output. Only the last output is
directly supervised. Intermediate outputs and hidden states remain attached to
the computation graph, allowing gradients to propagate through all $L$ steps.
Xavier-uniform initialization~\cite{pmlr-v9-glorot10a} is used for the input,
output and output-feedback matrices, orthogonal initialization~\cite{saxe2014exactsolutionsnonlineardynamics}
for hidden-feedback matrices, and zero initialization for biases.

\subsection{Loss and Optimization}
For a batch of standardized targets, let $M_{ij}$ be one when target $j$ of
example $i$ is observed and zero otherwise. The shared loss is masked mean
squared error:
\begin{equation}
 \mathcal{L}=\frac{\sum_{i=1}^{B}\sum_{j=1}^{m}
 M_{ij}(\hat z_{ij}-z_{ij})^2}{\sum_{i=1}^{B}\sum_{j=1}^{m}M_{ij}}.
\end{equation}
In code, observed entries are selected before subtraction to prevent masked
placeholders from contaminating the calculation. This gives each observed
standardized target entry equal weight.

Gradients are computed by backpropagation through time using PyTorch autograd.
Adam updates the weights with learning rate $0.001$ and its default moment
parameters. Batches contain up to 32 examples, and the total gradient norm is
clipped to 1.0 before each update. The search allows at most 100 epochs. There
is no dropout or weight decay. Examples are shuffled within the training
partition at each epoch, while the time order inside each window is preserved.
This does not randomly mix training and validation periods: the partitions remain
chronological, and recurrent context is reset between windows.

\subsection{Chronological Validation and Model Selection}
The last block of each dataset is reserved as a final test set before tuning.
The earlier development interval is evaluated using expanding-window
cross-validation: the training prefix grows, and each following validation block
is scored once. The step between folds equals the validation-block length.
Table~\ref{tab:split-settings} gives the actual row counts and input windows.
Only complete validation blocks are emitted; any development rows left after the
last complete block are still available for final training. Sliding-window
validation is not used in the reported experiments.

\begin{table*}[t]
\centering
\small
\caption{Preprocessing and temporal split settings. All sizes are row counts,
not numbers of constructed examples. The forecast horizon is one step
throughout.}
\label{tab:split-settings}
\begin{tabular}{lrrrrr}
\toprule
Dataset & Lookback & Initial training & Validation block & Reserved test & CV folds \\
\midrule
Air passengers & 12 & 72 & 12 & 24 & 4 \\
Minimum temperatures & 30 & 1461 & 365 & 730 & 4 \\
Smart home & 96 & 1344 & 288 & 672 & 7 \\
Sunspots & 132 & 2000 & 264 & 264 & 4 \\
AR(2) & 2 & 2000 & 500 & 1000 & 4 \\
\bottomrule
\end{tabular}
\end{table*}

Training duration is chosen inside each outer training fold, without consulting
its outer validation targets. The final 15\% of that training interval, rounded
up to whole rows, forms an inner early-stopping block. Separate scalers are
fitted on the inner fitting prefix. After each epoch, the inner stopping loss is
evaluated; an improvement must exceed $10^{-5}$ in standardized MSE. Training
ends after ten consecutive epochs without such improvement, or at 100 epochs,
and the epoch of the best qualifying improvement is retained.

The inner model is then discarded. A new model and optimizer are initialized
with the same seed, scalers are fitted on the entire outer training interval,
and the model is trained for the retained number of epochs. Only this refitted
model is scored on the outer validation block. This separates the choice of
training duration from the validation score used to compare hidden sizes.

For each architecture, hidden sizes $H\in\{8,16,32\}$ are compared using seed
42. A fold's selection score is the mean of its per-target root mean square
errors (RMSEs) in
training-standardized units; fold scores are then averaged with equal weight.
This avoids letting target units determine the multivariate ranking. It differs
from pooling all residuals before computing a single RMSE. Each architecture
retains its own lowest-scoring hidden size, with exact ties resolved in favour
of the smaller size. Across the five datasets, this produces 207 outer trials,
each including inner selection and outer refitting. Early stopping and the size
search constrain fitting, but do not guarantee the absence of underfitting or
overfitting.

\subsection{Final Evaluation and Reproducibility}
For each selected architecture and hidden size, the final epoch count is the
integer part of the median selected epoch across its cross-validation (CV)
folds, with a minimum of
one. New models are fitted on the complete development interval using seeds 11,
22 and 33 and development-only scalers. No test loss is used for early stopping
or model selection. This gives 45 final fits. Predictions are made one step
ahead using the latest available observations; the evaluation does not
recursively substitute model predictions for subsequently observed inputs.

For target $j$, errors in original units are summarized over the observed test
indices $\mathcal{O}_j$:
\begin{align}
 \mathrm{MAE}_j &= \frac{1}{|\mathcal{O}_j|}
     \sum_{i\in\mathcal{O}_j}|\hat y_{ij}-y_{ij}|,\\
 \mathrm{RMSE}_j &= \sqrt{\frac{1}{|\mathcal{O}_j|}
     \sum_{i\in\mathcal{O}_j}(\hat y_{ij}-y_{ij})^2}.
\end{align}
(MAE stands for mean absolute error.)
Persistence predicts each target's latest available value,
$\hat{\mathbf{y}}_{t+1}=\mathbf{y}_t$, and is evaluated on the same eligible
targets. Metrics are calculated separately for each final fit and reported as
their mean and sample standard deviation across seeds, rather than as the
performance of an ensemble. Both smart-home outputs are reported separately.

The runner saves configurations, data and implementation hashes, software
versions, fitted scalers, model checkpoints, epoch histories and timestamped
predictions. These artifacts support resumption and checking of experiment
consistency. The results-generation script recomputes errors from all saved
observed predictions and verifies identical evaluation rows across architectures
and seeds before creating report tables and figures. Automated checks also cover
temporal splitting, preprocessing leakage, masking, recurrent equations,
gradient propagation and independent window state. These checks establish
implementation consistency.
